\section{实践指南}

\subsection{推荐配置}

根据讲者的讨论和 DPM-Solver 的实践经验，以下是不同场景的推荐配置：

\begin{table}[H]
\centering
\begin{tabular}{lcc}
\toprule
\textbf{场景} & \textbf{推荐方法} & \textbf{步数} \\
\midrule
快速预览 / 交互式应用 & DPM-Solver++(2M) & 10$\sim$15 \\
高质量生成 & DPM-Solver++(2M) & 20$\sim$25 \\
最高质量（不计速度） & DPM-Solver++(3M) & 25$\sim$50 \\
\bottomrule
\end{tabular}
\caption{不同场景下的推荐 DPM-Solver 配置（2M = 二阶多步法，3M = 三阶多步法）}
\end{table}

\subsection{实现要点}

\begin{warningbox}{实现中的常见陷阱}
\begin{enumerate}
    \item \textbf{数值稳定性}：$e^{\lambda}$ 在 $\lambda$ 值较大时可能溢出，需要使用 log-space 计算
    \item \textbf{时间步方向}：扩散模型的"前向"是加噪（$t: 0 \to T$），"逆向"是去噪（$t: T \to 0$）。ODE 求解的方向是逆向的，$h = \lambda_t - \lambda_s$ 通常为\textbf{正值}（因为 $\lambda$ 随去噪递增）
    \item \textbf{多步法的预热}：二阶多步法在第一步没有历史缓存，需回退到一阶单步法
    \item \textbf{与 Classifier-free Guidance 的兼容性}：CFG 会改变有效的 ODE，DPM-Solver++ 对此做了专门的适配
\end{enumerate}
\end{warningbox}

\begin{knowledgebox}{DPM-Solver 在主流框架中的集成}
DPM-Solver/DPM-Solver++ 已被广泛集成在主流扩散模型库中：
\begin{itemize}
    \item \textbf{Diffusers}（Hugging Face）：\texttt{DPMSolverMultistepScheduler}
    \item \textbf{k-diffusion}：作为默认采样器之一
    \item \textbf{Stable Diffusion WebUI}：在 Sampler 下拉菜单中可选
\end{itemize}
在实际使用中，通常只需将 scheduler/sampler 切换为 DPM-Solver++，即可享受加速效果，无需修改模型本身。
\end{knowledgebox}

\subsection{本章小结}

DPM-Solver++ 的二阶多步法（2M）是实际应用中的最佳默认选择，兼顾速度和质量。实现时需注意数值稳定性和多步法预热等细节。

